An operator who controls field devices through a network must continue to give orders when
links have failures. But a device that obeys an isolated operator cannot know if the operator is
loyal. It is possible that an attacker has the key of the operator and cut the links. In the
2015 attack on the Ukrainian power grid, attackers opened breakers from the consoles of the
operators. We examine how command stops, and how command goes to a different person, when
the commander can be hostile and the attacker controls all links.

First, we show a limit. For each protocol, a hostile commander keeps a device for a time
not less than the outage that a loyal commander can continue through. The same
argument applies to a deputy who gets command on a timeout.

Second, we give a protocol, \sys{}. Command expires at a time that each device compares
with its clock, unless a quorum of officers signs a renewal. A specified number of loyal
officers always removes a commander, and hostile officers cannot remove a commander without
loyal officers. \sys{} is near the limit. The difference is not more than the sum of the
renewal interval, the collection time and six clock errors.

Third, we give the cost of \sys{} as formulas and compare it with four known designs. The
paper has the proofs. A model check of the rules and a simulation with a hostile commander
are not done at this time.
